\documentclass[a4paper, 12pt]{report}

\usepackage[spanish]{babel}
\usepackage[utf8]{inputenc}
\usepackage[left=4cm, right=4cm, top=4cm]{geometry}
\usepackage{textcomp}
\usepackage{booktabs}
\usepackage{amssymb}
\usepackage{bussproofs}
\usepackage{fancyhdr}
\usepackage{graphicx}
\usepackage{amsmath}
\usepackage{enumitem}
\usepackage{ifsym}
\usepackage{float}


\usepackage{hyperref}
\hypersetup{
    colorlinks=true,
    linkcolor=blue,
    filecolor=magenta,
    urlcolor=cyan,
}

\pagestyle{fancy}
\lhead{Almeida, Figueroa \& Ibarra}
\chead{Tarea 5}
\rhead{\today}

\begin{document}
\begin{titlepage}
    \centering
    {\scshape\Huge Universidad Nacional Autónoma de México \par}
    \vspace{1.25cm}
    {\scshape\huge Fundamentos de Bases de Datos\par}
    \vspace{1.25cm}
    {\huge\bfseries Tarea 5:\\ Dependencias y Normalización\par}
    \vspace{1.25cm}
    {\Large\textsc Almeida Rodríguez Jerónimo\par}
    \vspace{.1cm}
    {\large\texttt{418003815}\par}
    \vspace{0.25cm}
    {\Large\textsc Figueroa Sandoval Gerardo Emiliano\par}
    \vspace{.1cm}
    {\large\texttt{315241774}\par}
    \vspace{0.25cm}
    {\Large\textsc Ibarra Moreno Gisselle \par}
    \vspace{.1cm}
    {\large\texttt{315602193}\par}
    \vspace{1.5cm}
    \vfill
    \begin{figure}[hb!]
        \includegraphics[width=.3\textwidth]
            {../logos/escudo_f-ciencias.png}\hfill
        \includegraphics[width=.3\textwidth]
            {../logos/Escudo_UNAM.png}\hfill
    \end{figure}
\end{titlepage}

\begin{enumerate}
\item{
    \textbf{Preguntas de repaso:}
    \begin{itemize}
        \item 1
        \item 2
    \end{itemize}
}
\item {Proporciona algunos ejemplos que demuestren que las siguientes reglas 
	no son válidas:
\begin{enumerate}
	\item Si A → B, entonces B → A\\
	Sea la relación R(restaurante,ciudad,teléfono).\\
	Se tiene la DF $restaurante \rightarrow ciudad$, ya que un restaurante tiene 
	asociada una sola ciudad, pero notemos que la DF $ciudad \rightarrow 
	restaurante$ no se da, ya que una ciudad tiene varios restaurantes. \\
	Por lo que A $\rightarrow$ B no implica que B $\rightarrow$ A
\end{enumerate}}
\item {Para cada uno de los esquemas que se muestran a continuación:
\begin{enumerate}
	\item R(A,B,C,D,E) con F = \{AB $\rightarrow$ CD, E $\rightarrow$ C, D
		$\rightarrow$ B\}
	\item R(A,B,C,D,E)con F = \{AB $\rightarrow$ C, DE $\rightarrow$ C, B $\rightarrow$ D\}
\end{enumerate}
\begin{itemize}
	\item Especifica de  ser  posible dos DF  no  triviales que  se  pueden  
	derivar  de  las  dependencias funcionales dadas.\\
	a. \{AB $\rightarrow$ C, AB $\rightarrow$ D\}.\\
	\item Indica alguna llave candidata para R.\\
	a.\{ABE\}+=\{ABCDE\}\\
	\item Especifica todas las violaciones a la BCNF\\
	a.\{AB $\rightarrow$ CD, E $\rightarrow$ C, D
	$\rightarrow$ B\}
	\item Normaliza de acuerdo a BCNF, asegúrate de indicar cuáles son las 
	relaciones resultantes con sus respectivas dependencias funcionales:\\
	a. Tomamos la violación AB $\rightarrow$ CD.\\
	Obtenemos las relación S(A,B,C,D) con dependecias \{AB $\rightarrow$ CD,
	D $\rightarrow$ B\} y la relación T(A,B,E) con dependencias \{ABE $\rightarrow$ ABE\}.\\
	\{AB\}+=\{ABCD\} es una llave para S, entonces tomamos la violación D $\rightarrow$ B.\\
	Obtenemos la relación U(D,B) con dependencia \{D $\rightarrow$ B\} y la
	relación V(D,A,C) con dependencia \{DAC $\rightarrow$ DAC\}.\\
	Por lo tanto el esquema en BCNF es U(D,B), V(D,A,C) y T(A,B,E).
\end{itemize}}
\item Para cada una de  las  siguientes  relaciones  con  su  respectivo  conjunto  de  dependencias funcionales:
\begin{enumerate}
	\item R(A,B,C,D,E,F)con F = \{B $\rightarrow$ D, B $\rightarrow$ E, D $\rightarrow$ F, AB $\rightarrow$ C\}
	\item R(A,B,C,D,E)con F = \{A $\rightarrow$ BC, B $\rightarrow$ D, CD$\rightarrow$ E, E $\rightarrow$ A \}
\end{enumerate}
\begin{itemize}
	\item Indica todas las violaciones a la 3NF\\
	b.\{A\}+ = \{ABCDE\} y \{E\}+ = \{EABCD\} son llaves, entonces la única 
	violación a la 3NF es B $\rightarrow$ D.
	\item Normaliza de acuerdo a la 3NF\\
	b. Superfluos por la izquiera:  CD $\rightarrow$ E\\
	¿C es superfluo? D $\rightarrow$ E, \{D\}+=\{D\} entonces C no es superfluo.\\
	¿D es superfluo? C $\rightarrow$ E, \{C\}+=\{C\} entonces D no es superfluo.\\
	Superfluos por la derecha: A $\rightarrow$ BC\\
	¿B es superfluo? A $\rightarrow$ C, F'=\{A$\rightarrow$C, B$\rightarrow$D, CD$\rightarrow$E, E$\rightarrow$A\}.\\
	\{A\}+= \{AC\} por lo tanto, B no es superfluo.\\
	¿C es superfluo? A$\rightarrow$B F''=\{A$\rightarrow$B, B$\rightarrow$D, CD$\rightarrow$E, E$\rightarrow$A\}.\\
	\{A\}+= \{AB\} por lo tanto, C no es superfluo.\\
	Entonces F ya es el mínimo.\\
	S(A,B,C), T(B,D), U(C,D,E) y V(E,A) es el esquema en 3NF.
\end{itemize}
\item 

\item Para  cada  uno  de  los  esquemas,  con  su  respectivo  conjunto  de  dependencias  multivaluadas, resuelve los siguientes puntos:\\
\begin{enumerate}
	\item R(A,B,C,D)con DMV = {AB $\twoheadrightarrow$C, B $\rightarrow$ D}
	\item R(A,B,C,D,E)con DMV = { A $\twoheadrightarrow$B, AB $\rightarrow$ C, A $\rightarrow$ D, AB $\rightarrow$ E}\\
\end{enumerate}
\begin{itemize}
	\item Encuentra todas las violaciones a la 4NF.
	\{AB\}+=\{ABCDE\} es una llave. Las violaciones son \{A$\twoheadrightarrow$B\}.\\
	\item Normaliza de acuerdo a la 4NF.\\
	Tomamos la violación A$\twoheadrightarrow$B.\\
	Obtenemos las relaciones S(A,B) con DMV=\{A$\twoheadrightarrow$B\} y
	T(A,C,D,E) con DMV=\{A$\rightarrow$\}.\\
	\{ACE\}+=\{ACED\} es una llave para T, entonces tomamos la violación
	A$\rightarrow$D.\\
	Obtenemos las relaciones U(A,D) con DMV=\{A$\rightarrow$\} y V(A,C,E) con
	DMV=\{ACE$\rightarrow$\}.\\
	Por lo tanto, el esquema en 4NF es S(A,B), U(A,D) y V(A,C,E).
\end{itemize}
\end{enumerate}

\end{document}
